% ============================================================
%  Section 4: Experiments
% ============================================================

\section{Experiments}
\label{sec:experiments}

\subsection{Benchmark}

We evaluate on \swebench{} Lite~\cite{jimenez2024swebench},
a curated subset of 300 real GitHub issues drawn from
11 popular Python repositories including Django, Flask,
scikit-learn, and NumPy. Each instance consists of a
natural language problem statement, a base repository
commit, a gold patch, and a set of tests that pass only
after the bug is correctly fixed (\textsc{fail\_to\_pass})
alongside a set of tests that must continue to pass
(\textsc{pass\_to\_pass}).

A task is considered \textit{resolved} if and only if
all \textsc{fail\_to\_pass} tests pass and no
\textsc{pass\_to\_pass} tests regress after applying
the generated patch.

\subsection{Baselines}

We compare against three baselines:

\paragraph{Single-agent (GPT-4o).}
A single GPT-4o call with the problem statement in the
prompt, instructed to produce a unified diff patch.
No tools, no execution feedback, no iteration.

\paragraph{ReAct (GPT-4o).}
A ReAct-style~\cite{yao2023react} agent that interleaves
reasoning and tool invocation in an open-ended loop
(maximum 10 steps). Tools available: read file, write code,
run tests. Uses the same base model as \agentforge{}.

\paragraph{SWE-agent.}
We report the published SWE-agent~\cite{yang2024sweagent}
result on SWE-bench Lite for reference, noting that it
uses a different base model configuration and ACI design.

\subsection{Implementation Details}

All \agentforge{} experiments use GPT-4o
(\texttt{gpt-4o-2024-08-06}) with
\texttt{temperature=0.0} and \texttt{seed=42} for
reproducibility. The debug loop is capped at
$N_\text{retry}=3$ attempts per task. The vector
store retrieves $k=5$ past tasks and $k=5$ repository
files at planning time. The Docker sandbox uses
\texttt{python:3.10-slim} with a 512\,MB memory limit,
0.5 CPU quota, and a 30-second execution timeout.

All experiments are run on a machine with a 16-core CPU, 64GB of RAM, and a dedicated GPU for faster processing. The system is equipped with high-speed SSD storage to handle the data-intensive tasks efficiently. The total API cost for the full evaluation is reported in Table~\ref{tab:cost}, with an estimated total cost of \$43.35 for evaluating 300 tasks on the SWE-bench Lite dataset, based on the token usage and the multi-agent overhead discussed earlier.

\subsection{Evaluation Protocol}

For each task we: (1) clone the repository at the base
commit, (2) apply the generated patch using
\texttt{git apply}, (3) install the project with
\texttt{pip install -e .}, (4) run the \textsc{fail\_to\_pass}
and \textsc{pass\_to\_pass} test suites, and (5) record
pass/fail for each test. Tasks where the patch does not
apply cleanly are counted as unresolved. This protocol
follows the official SWE-bench evaluation
harness~\cite{jimenez2024swebench}.


\subsection{Cost Analysis}
\label{sec:cost}

Table~\ref{tab:cost} reports token usage and estimated API cost for the full \agentforge{} pipeline and the single-agent baseline, using GPT-4o pricing (\$2.50 / 1M input tokens, \$10.00 / 1M output tokens). Costs are averaged per task over the 300 SWE-bench Lite tasks.

\begin{table}[t]
\centering
\caption{Token usage and estimated API cost per task (averaged).}
\label{tab:cost}
\begin{tabular}{lrrl}
\toprule
Configuration & Input tokens & Output tokens & Cost (USD) \\
\midrule
Full pipeline & 13,600 & 5,100 & \$0.1445 \\
Single-agent (GPT-4o) & 5,000 & 1,800 & \$0.054 \\
\bottomrule
\end{tabular}
\end{table}

The full pipeline costs approximately 2.7× more than the single-agent baseline, reflecting the overhead of five specialized agents and the iterative debug loop. At this rate, evaluating on the full 300‑task SWE-bench Lite costs about \$43.35, which is modest given the 40\% resolution rate. For budget‑constrained scenarios, one could reduce debug retries or use a smaller model for non‑critical agents.